%% Font size %%
\documentclass[11pt]{article}

%% Load the custom package
\usepackage{Mathdoc}

%% Numéro de séquence %% Titre de la séquence %%
\renewcommand{\centerhead}{S14- Probabilités - Exercices Divers}

%% Spacing commands %%
\renewcommand{\baselinestretch}{1} \setlength{\parindent}{0pt}

%%\pagestyle{empty}

\begin{document}

\begin{exercice}[1][Reconnaître une expérience aléatoire]

Pour chaque situation, dire s'il s'agit ou non d'une expérience aléatoire.

\begin{enumerate}
    \item Lancer un dé à 6 faces.
    \item Mesurer la longueur d'une table avec une règle.
    \item Tirer une carte dans un jeu.
    \item Choisir au hasard une boule dans un sac.
    \item Calculer $15+7$.
\end{enumerate}

\end{exercice}


\begin{exercice}[1][Décrire les issues possibles]

On lance une pièce de monnaie.

\begin{enumerate}
    \item Donner toutes les issues possibles.
    \item Combien y a-t-il d'issues possibles ?
\end{enumerate}

\end{exercice}


\begin{exercice}[1][Issues d'un dé]

On lance un dé à 6 faces numérotées de 1 à 6.

\begin{enumerate}
    \item Donner toutes les issues possibles.
    \item Donner l'événement :
    \begin{enumerate}
        \item « obtenir un nombre pair » ;
        \item « obtenir un nombre supérieur à 4 ».
    \end{enumerate}
\end{enumerate}

\end{exercice}


\begin{exercice}[1][Événements]

On tire une boule dans un sac contenant des boules rouges, bleues et vertes.

Associer chaque événement à ses issues possibles.

\begin{enumerate}
    \item « Obtenir une boule rouge »
    \item « Obtenir une boule qui n'est pas verte »
    \item « Obtenir une boule bleue ou verte »
\end{enumerate}

\end{exercice}
\end{document}
